\scanpage{20}
% L1–3
Soit $D$ une géométrie de drapeaux quasi-simpliciale
\[
\deleted{D:\bigl(\unsure{A}\big/\underline{P_f}^{*}(\mathbb N)\bigr)^{\circ}\longrightarrow(\mathrm{Ens})}
\]
\deleted{une \uline{\textquotedblleft géométrie de drapeaux avec fonction-dimension\textquotedblright},}
% L4–5
[i.e avec \unsure{$D$\,:|$D$\,;|$D$} un syst. d'ordres totaux compatibles, sur les simplexes de $\unsure{\mathscr D|\mathscr F|\Phi}=D_0$ (ens des él. minimaux de $\mathscr D$)]
% L6–7
d'où
\[
X=|D|
\]
% L8–10
espace topologique, muni de parties fermées (en fait, des simplexes)
\[
X(d)\qquad d\in D
\]
% L11–13
définissant une \deleted{\unsure{structure}} \inserted{stratification} \unsure{cônique|conique|canonique} de $X$ par des simplexes (les $X^{\circ}(d)=X(d)\setminus\bigcup_{d'<d}X(d')$ \unsure{font|forment} une partition de $X$ par des simplexes ouverts...).

% L14
Pour tout $F\in\mathscr D$, soit
% L15–17
\[
\left\{\begin{aligned}
|F|&=\bigcup_{\substack{d\in D_{\unsure{*|\bullet}}\\
\text{t.q. }\deleted{\unsure{\sup I_d\in F}}\\ \sup(I_d)=F}}X(d)\\
|F|^{\circ}&=\bigcup_{\text{idem}}X(d)^{\circ}
\end{aligned}\right.
\]
% L15–17 right
\margin{$\displaystyle\bigcup_{F'<F}|F'|^{\circ}$\\
$\parallel$\\
$\displaystyle |F|-|F|^{\circ}\mathrel{\overset{?}{=}}\bigcup_{\substack{\deleted{\unsure{d\in D_*}}\\F'<F}}|F'|$\\
$\parallel$\\
$\displaystyle\bigcup_{\substack{d\in D_*\\\sup I_d<F}}X(d)$}
% L18–20
\margin{?}
Quand peut-on dire que la famille \inserted{des} $(|F|)_{F\in\mathscr D}$
\deleted{\unsure{forme}} \inserted{est celle des cellules fermées} d'une \uline{stratification} de $X$, dont les cellules ouvertes sont les $|F|^{\circ}$ ?

% L21
\margin{En fait}
\deleted{a) $|F|$ fermés $\bigcup_{F\in\mathscr D}|F|=\bigcup_{d\in D_*}X(d)=X$}
% L22–24
a) Les $|F|^{\circ}$ \inserted{loc. fermés ??} font partition de $X$
[\uline{NB} $|F|^{\circ}\ni\deleted{\unsure{\varnothing}}\unsure{\mathrm{V}|X}(F)$
\inserted{\unsure{espace ambiant|espaces ambiants|le point représentant|les sommets maximaux}} donc $|F|^{\circ}\ne\varnothing$]

% L25
b) $|F|=\overline{|F|^{\circ}}$ (car $X(d)=\overline{X(d)^{\circ}}$) \margin{réunion de \textquotedblleft cellules ouvertes\textquotedblright\ \unsure{...}}

% L26
\deleted{c) $|F|^{\circ}=|F|\setminus\bigcup_{\unsure{G< F}}\bigcup_{\substack{d\in D_*\\\sup(I_d)=F}}(X(d)$}
% archive sheet number: 19
